%% kapa: exam=1 round=37 year=2026 period=2 form=A subject=appraisal_law subject_ko=감정평가관계법규 seq=40 number=40
%% source: Q-Net 감정평가사 기출문제 (cst00309) artlSeq=5258924 file=raw/1008/5258924/2026년 제37회 감정평가사 1차 2교시.pdf
%% conversion: kordoc 4.24.51 -> tools/gen1.py (auto); review=proofread-sonnet; flags=none
\begin{question}{40}
동산ㆍ채권 등의 담보에 관한 법령상 담보등기에 관한 설명으로 옳지 않은 것은?
\begin{choices}
  \item ｢부가가치세법｣에 따라 사업자등록을 한 사람이 담보약정에 따라 금전의 지급을 목적으로 하는 지명채권을 담보로 제공하는 경우에는 담보등기를 할 수 있다.
  \item 포괄승계로 인한 등기는 등기권리자 단독으로 신청할 수 있다.
  \item 등기관이 등기를 마친 경우 그 등기는 접수한 때부터 효력을 발생한다.
  \item 담보등기는 동산담보권, 채권담보권 또는 지식재산권담보권의 설정, 이전, 변경, 말소 또는 연장에 대하여 한다.
  \item 등기관이 최초 담보권설정등기를 마쳤을 때에는 등기필정보를 등기권리자와 담보권설정자에게 통지하여야 한다.
\end{choices}
\end{question}
\answer{4}
